\documentclass{article}
\usepackage[utf8]{inputenc}
\usepackage[russian]{babel}
\usepackage{amsmath, amssymb}
\usepackage{enumitem}

\title{Вопросы на защиту}
\author{}
\date{}

\begin{document}

\maketitle

\begin{enumerate}
\item Сформулируйте задачу Коши для обыкновенного дифференциального уравнения $n$-го порядка. 

\item Можно ли решить уравнение из задания~(2) методом, используемым в задании~(3)? Обоснуйте ответ, указав ограничения применимости метода неопределённых коэффициентов.

\item Можно ли решить уравнение из задания~(3) методом вариации произвольных постоянных, применяемым в задании~(2)? Почему этот метод применим или нецелесообразен в данном случае?

\item Возможно ли решить уравнение из задания~(4) методом, используемым в заданиях~(3) или~(2)? Объясните, какие свойства коэффициентов уравнения препятствуют или допускают применение этого метода.

\item Запишите общий вид специальной правой части линейного дифференциального уравнения с постоянными коэффициентами.

\item Для системы из задания~(5) выпишите соответствующую однородную систему и приведите её к одному линейному дифференциальному уравнению второго порядка (указать, какую переменную исключаете).

\item Какой объект в теории линейных дифференциальных уравнений второго порядка с постоянными коэффициентами является аналогом собственных значений матрицы в задаче на нахождение матричной экспоненты?

\item В чём заключается идея понижения порядка дифференциального уравнения? При каких условиях этот приём применим и какие замены при этом используются?

\item Какую информацию о пространстве решений линейного однородного уравнения формула Остроградского--Лиувилля позволяет получить?

\item Опишите общий алгоритм метода вариации произвольных постоянных для линейных неоднородных дифференциальных уравнений.
\end{enumerate}

% \newpage
% \begin{enumerate}
% \item \textbf{Задача Коши для ОДУ $n$-го порядка.}  
% Для уравнения
% \[
% y^{(n)} = f(x, y, y', \dots, y^{(n-1)})
% \]
% задача Коши состоит в нахождении решения, удовлетворяющего условиям
% \[
% y(x_0)=y_0,\; y'(x_0)=y_1,\; \dots,\; y^{(n-1)}(x_0)=y_{n-1}.
% \]

% \item \textbf{Применимость метода неопределённых коэффициентов к заданию (2).}  
% Нет, нельзя. Метод неопределённых коэффициентов применим только к линейным уравнениям с постоянными коэффициентами и специальной правой частью (полином, экспонента, синус/косинус и их произведения). В задании~(2) правая часть содержит множитель $t^{3/2}$, не являющийся допустимым.

% \item \textbf{Применимость метода вариации постоянных к заданию (3).}  
% Да, можно. Метод вариации произвольных постоянных применим к любым линейным неоднородным уравнениям, в том числе с постоянными коэффициентами и произвольной правой частью. Однако он вычислительно сложнее, чем метод неопределённых коэффициентов.

% \item \textbf{Применимость методов из (2) и (3) к заданию (4).}  
% Нет. Уравнение из задания~(4) имеет переменные коэффициенты и является уравнением Эйлера--Коши с логарифмическим множителем, поэтому методы для уравнений с постоянными коэффициентами неприменимы.

% \item \textbf{Специальная правая часть для метода неопределённых коэффициентов.}  
% Общий вид:
% \[
% f(x)=e^{\alpha x}\bigl(P_n(x)\cos \beta x + Q_m(x)\sin \beta x\bigr),
% \]
% где $P_n, Q_m$ — многочлены, $\alpha,\beta\in\mathbb{R}$ — параметры.

% \item \textbf{Однородная система и приведение к ОДУ второго порядка (задание 5).}  
% Однородная система:
% \[
% \begin{cases}
% \dot x_1 = x_1 - x_2,\\
% \dot x_2 = 2x_1 - x_2.
% \end{cases}
% \]
% Исключая $x_2$, получаем для $x_1$:
% \[
% x_1'' - 2x_1' + x_1 = 0.
% \]

% \item \textbf{Аналог собственных значений матрицы.}  
% Корни характеристического уравнения линейного дифференциального уравнения с постоянными коэффициентами являются аналогом собственных значений матрицы.

% \item \textbf{Идея понижения порядка.}  
% Понижение порядка основано на введении новой переменной (например, $p=y'$ или $p(y)=y'$), что позволяет свести уравнение второго порядка к уравнению первого порядка. Метод применим при отсутствии явной зависимости от независимой переменной или функции.

% \item \textbf{Формула Остроградского--Лиувилля.}  
% Формула позволяет найти вронскиан фундаментальной системы решений и установить линейную независимость решений линейного однородного уравнения.

% \item \textbf{Метод вариации произвольных постоянных.}  
% Алгоритм:
% \begin{enumerate}
% \item Найти фундаментальную систему решений однородного уравнения.
% \item Заменить постоянные на функции.
% \item Получить систему для этих функций.
% \item Проинтегрировать и подставить в общее решение.
% \end{enumerate}
% \end{enumerate}

% \bigskip
% \begin{center}
% \begin{tabular}{|c|p{10cm}|}
% \hline
% \textbf{№} & \textbf{Краткий ответ} \\
% \hline
% 1 & Значения функции и её производных до порядка $n-1$ в точке $x_0$ \\
% 2 & Нет, правая часть не специального вида \\
% 3 & Да, метод универсален, но громоздок \\
% 4 & Нет, коэффициенты переменные \\
% 5 & Экспоненты, полиномы, синусы и косинусы \\
% 6 & Получается уравнение $x_1''-2x_1'+x_1=0$ \\
% 7 & Корни характеристического уравнения \\
% 8 & Замена $y'$ или использование отсутствия аргумента \\
% 9 & Линейная независимость решений \\
% 10 & Замена констант функциями и интегрирование \\
% \hline
% \end{tabular}
% \end{center}


\end{document}